\scanpage{46}\mspage{33}

composé de celui-ci et de $(\unsure{*|\alpha})$ \deleted{\unsure{étant}} (base du diagramme \unsure{polyédral|polyèdral} $(**)$) étant \unsure{cocartésiens|cocartésien}, il en est de $=$ de $(*)$.

\uline{Lemme} 2.7. Soit $(I_\alpha)_{\alpha\le\alpha_0}$ une suite ordinale croissante d'ouverts de $I$, \deleted{dont réunion $I_{\alpha_0}$. Alors, \unsure{et|de}} Alors $u_{I_0,I_{\alpha_0}}$$:$\inserted{$Y_{I_0,I_{\alpha_0}}\unsure{\mathbin{-}|\longrightarrow}Y_{I_{\alpha_0}}$} est composée transfinie, indexée par l'ens des ordinaux $\deleted{\unsure{\alpha<\alpha_0}}\inserted{-1\le}\alpha<\alpha_0$, d'une suite de flèches $v_\alpha$, où chaque $v_\alpha$ est déduit de $u_{I_\alpha,I_{\alpha+1}}$ par cochangement de base (et de $u_{I_0}:X_{I_0}\longrightarrow Y_{I_0}$ pour $\alpha=0$).

Pour le voir, on introduit la suite transfinie d'objets $Z_\alpha$\inserted{$=Y_{I_\alpha,I_{\alpha_0}}$}, définis par le \deleted{$Z_\alpha=$} diagramme

\diagram{\begin{tikzcd}[ampersand replacement=\&,row sep=large,column sep=large]
 \deleted{\unsure{Z|X}}X_{I_\alpha}\arrow[r]\arrow[d] \& X_{I_{\alpha_0}}\arrow[d,dash]\\
 Y_{I_\alpha}\arrow[r] \&\deleted{\unsure{Z|Y}}\unsure{Y|Z}_{I_\alpha,I_{\alpha_0}}=Z_\alpha\arrow[d]\\
 \&Y_{I_{\alpha_0}},
\end{tikzcd}}

on a \deleted{\unsure{du}}

\[
 Z_0=Y_{I_0,I_{\alpha_0}},\quad Z_{\alpha_0}=Y_{I_{\deleted{0}\alpha_0},I_{\alpha_0}},\quad
 Z_\alpha=\underset{\beta<\alpha}{\unsure{\underrightarrow{\lim}|\underline{\lim}}}\,Z_\beta
 \quad\text{si }\unsure{\alpha}
\]

% archive sheet number: 45

